\documentclass{article}
\usepackage{amsthm}
\usepackage{amssymb}
\usepackage{url}
\usepackage{imprimatur}
\newtheorem{theorem}{Theorem}
% The paper's environments, as imprimatur.toml's [aliases] declares them.
\leanblockalias{leangrammar}{}
\leanblockalias{leantheorem}{theorem}
\leanblockalias{leanconventions}{}
\leanblockalias{leannotation}{}

\title{imprimatur: papers that say what Lean proved}
\author{The imprimatur authors}
\date{}

\begin{document}
\maketitle

\begin{abstract}
A paper about a Lean development restates its definitions and theorems in
mathematical notation. imprimatur generates those restatements from the
development, keeps the text the authors wrote around them, and checks that
the paper still says what Lean proved. We state what that check guarantees,
and prove it for a model of imprimatur's splicer in Lean. This paper is
itself built with imprimatur: every displayed definition and theorem below
is generated from that Lean development, and every hand-written claim about
it is checked against it.
\end{abstract}

\section{Introduction}

A paper and its Lean development drift apart. A theorem is weakened in Lean
and not in the paper; a definition gains a case; a lemma the paper cites is
proved with \texttt{sorry}. imprimatur closes the gap in two ways. Blocks of
the paper are \emph{generated}: the extractor renders a Lean constant's
statement in the paper's notation, and the splicer writes it into the paper,
keeping any text the authors wrote in its \emph{slots}. Hand-written prose
about a constant is \emph{attested}: a \verb|\leancheck| beside it records the
hash of the statement the prose was checked against, and the check fails if
the statement changes.

What does it mean for this to be correct? A paper is correct with respect to
its development when what it shows is what Lean proved. We state that as two
theorems. The first, about the splicer, says what a successful
\texttt{imprimatur check} guarantees; we prove it for a model of the
splicer (Section~\ref{sec:check}). The second, about the notation, says that
a rendering can be read back as the statement it came from; imprimatur
checks it for each formula it renders (Section~\ref{sec:faithful}).

\section{Documents and bundles}

A document is a list of pieces: hand-written text, generated blocks, and
checks on hand-written prose. A block's body is a list of segments: text
generated from Lean, and slots, whose text the authors own.

\begin{leangrammar}[sum=g4GdhZSliS]{SpliceModel.Seg}{Segments}
\label{gram:SpliceModel.Seg}\begin{LeanGrammar}
\LeanProduction{\sigma}{A segment of a generated block.}{\mathsf{gen}(s)}{\leanslot{SpliceModel.Seg.gen}{text generated from Lean}}
\LeanAlternative{\mathsf{slot}_{i}(t)}{\leanslot{SpliceModel.Seg.slot}{a slot, whose text the authors own}}
\end{LeanGrammar}
\end{leangrammar}

\begin{leangrammar}[sum=iQOwawJZY6]{SpliceModel.Piece}{Pieces}
\label{gram:SpliceModel.Piece}\begin{LeanGrammar}
\LeanProduction{p}{A piece of a document.}{\mathsf{text}(s)}{\leanslot{SpliceModel.Piece.text}{hand-written text}}
\LeanAlternative{\mathsf{block}_{k}(\mathit{body})}{\leanslot{SpliceModel.Piece.block}{a generated block for a key}}
\LeanAlternative{\mathsf{checked}_{k}(h)}{\leanslot{SpliceModel.Piece.check}{hand-written prose, checked against hash $h$}}
\end{LeanGrammar}
\end{leangrammar}

A \emph{bundle}, exported from a Lean environment, gives each key the
rendering of its statement, the statement's hash, and whether its proof is
trusted, that is, uses no \texttt{sorry} and only the permitted
axioms\leancheck[stmt=81ba55cf6570, sig=401766da2f98]{SpliceModel.Entry, SpliceModel.extract}.

Regenerating a block fills the rendering's slots with the text the authors
wrote in the old block, slot by slot, and with the rendering's own text for a
slot they have not written:

\begin{center}
\begin{leanblock}[view=math, def, sum=BhwOfgB0yt]{SpliceModel.fill}{Filling slots}
\LeanMath{\begin{aligned}\mathsf{fill}(\mathit{old}, []) &= [] \\ \mathsf{fill}(\mathit{old}, \mathsf{gen}(s) \mathbin{::} r) &= \mathsf{gen}(s) \mathbin{::} \mathsf{fill}(\mathit{old}, r) \\ \mathsf{fill}(\mathit{old}, \mathsf{slot}_{i}(t) \mathbin{::} r) &= \mathsf{slot}_{i}(\mathit{old}[i] \mathbin{\mathsf{else}} t) \mathbin{::} \mathsf{fill}(\mathit{old}, r)\end{aligned}}
\end{leanblock}
\end{center}

\verb|imprimatur apply| regenerates every block whose key has an entry,
and leaves every other piece as it
is\leancheck[stmt=6deca83854c0, sig=7546958bfb76]{SpliceModel.apply, SpliceModel.regen}. \verb|imprimatur check|
reports what is not as the bundle says:

\begin{leangrammar}[sum=5CDH8zS1mO]{SpliceModel.Diag}{Diagnostics}
\label{gram:SpliceModel.Diag}\begin{LeanGrammar}
\LeanProduction{x}{What \LeanCodeInline{check} reports.}{\mathsf{broken}(k)}{\leanslot{SpliceModel.Diag.brokenRef}{the key has no entry}}
\LeanAlternative{\mathsf{untrusted}(k)}{\leanslot{SpliceModel.Diag.untrusted}{the key's proof is not trusted}}
\LeanAlternative{\mathsf{stale}(k)}{\leanslot{SpliceModel.Diag.stale}{the block is not the key's rendering}}
\LeanAlternative{\mathsf{drift}(k)}{\leanslot{SpliceModel.Diag.drift}{the statement has changed since the prose was checked}}
\end{LeanGrammar}
\end{leangrammar}

A key without an entry is a broken reference; an entry whose proof is not
trusted is reported; a block that is not the rendering, with its slots
filled, is stale; and a check whose recorded hash is not the statement's is
drift\leancheck[stmt=40c927a974cb, sig=680ab0e4d2e0]{SpliceModel.diags, SpliceModel.check}.

\section{The splicer's laws}

Four laws make \verb|apply| safe to run on a paper the authors are editing.
First, it changes nothing they wrote outside generated blocks. Write
$\lfloor d \rfloor$ for the outline of document $d$: its text, its checks,
and its blocks without their bodies.

\begin{leantheorem}[sum=3p5XLvWAci]{Locality}{Locality}
\label{thm:Locality}
\LeanProse{Let $b \in \mathit{Bundle}$ and $d \in \mathit{Doc}$. Then $\lfloor \mathsf{apply}_{b}(d) \rfloor = \lfloor d \rfloor$.}
\end{leantheorem}

Second, running it again changes nothing more. This needs each rendering's
slots to be distinct, which we write $\mathsf{wf}(b)$\leancheck[stmt=9891fff10987, sig=9d550ab02488]{SpliceModel.WF}.

\begin{leantheorem}[sum=rUqQQUgoM3]{Idempotence}{Idempotence}
\label{thm:Idempotence}
\LeanProse{Let $b \in \mathit{Bundle}$ and $d \in \mathit{Doc}$. Assume $\mathsf{wf}(b)$. Then $\mathsf{apply}_{b}(\mathsf{apply}_{b}(d)) = \mathsf{apply}_{b}(d)$.}
\end{leantheorem}

Third, \verb|check| agrees with \verb|apply|: after it, no block is stale.

\begin{leantheorem}[sum=ajKQ-ezoOC]{Fixed point}{Fixed point}
\label{thm:Fixed-point}
\LeanProse{Let $b \in \mathit{Bundle}$, $d \in \mathit{Doc}$ and $k \in \mathit{Key}$. Assume $\mathsf{wf}(b)$. Then $\mathsf{stale}(k) \notin \mathsf{check}_{b}(\mathsf{apply}_{b}(d))$.}
\end{leantheorem}

Fourth, the text the authors wrote in a slot survives regeneration, as long
as the rendering still has that slot.

\begin{leantheorem}[sum=4WbuqUPxi1]{Slot preservation}{Slot preservation}
\label{thm:Slot-preservation}
\LeanProse{Let $i \in \mathit{SlotId}$, $t \in \mathit{String}$, $\mathit{old} \in \mathsf{List}(\mathit{SlotId} \times \mathit{String})$ and $r \in \mathit{Body}$. Assume $i \in \mathsf{ids}(r)$ and $\mathit{old}[i] = t$. Then $\mathsf{slots}(\mathsf{fill}(\mathit{old}, r))[i] = t$.}
\end{leantheorem}

Accepting the current hashes, \verb|imprimatur accept|, likewise leaves no
drift for \verb|check| to report\leancheck[stmt=ac37d7be7013, sig=2d6e269c5e16]{SpliceModel.not_drift_after_accept}.

\section{What \texttt{check} guarantees}
\label{sec:check}

An environment, for a paper, is the statement of each constant and whether
its proof is trusted. The bundle the extractor exports from it renders and
hashes each statement. If \verb|check| reports nothing against that bundle,
the paper shows what Lean proved: every block is the rendering of a trusted
constant's statement, with the authors' slot text in it,

\begin{leantheorem}[sum=xC0mGigpH_]{Soundness for blocks}{Soundness for blocks}
\label{thm:Soundness-for-blocks}
\LeanProse{Let $\mathit{Stmt}$ be a type, $\mathit{render} \in \mathit{Stmt} \to \mathit{Body}$, $\mathit{hash} \in \mathit{Stmt} \to \mathbb{N}$, $E \in \mathit{Env}(\mathit{Stmt})$, $d \in \mathit{Doc}$, $k \in \mathit{Key}$ and $\mathit{body} \in \mathit{Body}$. Assume $\mathsf{check}_{\mathsf{export}(E)}(d) = []$ and $\mathsf{block}_{k}(\mathit{body}) \in d$. Then there is $T$ such that (1)~$\mathit{stmt}_{E}(k) = T$, (2)~$\mathit{trusted}_{E}(k)$ and (3)~$\mathit{body} = \mathsf{fill}(\mathsf{slots}(\mathit{body}), \mathit{render}(T))$.}
\end{leantheorem}

and every hand-written claim faces a trusted constant whose statement has
the hash recorded when the claim was checked.

\begin{leantheorem}[sum=V_NF2sG0U7]{Soundness for checked prose}{Soundness for checked prose}
\label{thm:Soundness-for-checked-prose}
\LeanProse{Let $\mathit{Stmt}$ be a type, $\mathit{render} \in \mathit{Stmt} \to \mathit{Body}$, $\mathit{hash} \in \mathit{Stmt} \to \mathbb{N}$, $E \in \mathit{Env}(\mathit{Stmt})$, $d \in \mathit{Doc}$, $k \in \mathit{Key}$ and $h \in \mathbb{N}$. Assume $\mathsf{check}_{\mathsf{export}(E)}(d) = []$ and $\mathsf{checked}_{k}(h) \in d$. Then there is $T$ such that (1)~$\mathit{stmt}_{E}(k) = T$, (2)~$\mathit{trusted}_{E}(k)$ and (3)~$\mathit{hash}(T) = h$.}
\end{leantheorem}

The hash stands for the statement only when hashes do not collide. Assuming
that, prose checked against a statement still faces that statement.

\begin{leantheorem}[sum=nUracosHsM]{Attested prose}{Attested prose}
\label{thm:Attested-prose}
\LeanProse{Let $\mathit{Stmt}$ be a type, $\mathit{render} \in \mathit{Stmt} \to \mathit{Body}$, $\mathit{hash} \in \mathit{Stmt} \to \mathbb{N}$, $E \in \mathit{Env}(\mathit{Stmt})$, $d \in \mathit{Doc}$, $k \in \mathit{Key}$ and $T_{0} \in \mathit{Stmt}$. Assume $\forall T : \mathit{Stmt},\ \forall T' : \mathit{Stmt},\ \mathit{hash}(T) = \mathit{hash}(T') \implies T = T'$, $\mathsf{check}_{\mathsf{export}(E)}(d) = []$ and $\mathsf{checked}_{k}(\mathit{hash}(T_{0})) \in d$. Then $\mathit{stmt}_{E}(k) = T_{0}$.}
\end{leantheorem}

The splicer that runs on real papers works on \TeX{} source, not on
documents already parsed into pieces. Its scanner and its edits are what the
model leaves out, and they are property-tested against the same four laws.

\section{Rendering is faithful}
\label{sec:faithful}

Soundness says the text is the \emph{rendering} of the statement. A reader
must also recover the statement from it: this is Pollack's question of
believing a machine-checked proof~\cite{Pollack98}, and Wiedijk's
\emph{Pollack-consistency}, that printing a term and parsing it back gives
the same term~\cite{Wiedijk12}.

The notation module defines both directions. Read as a grammar, each of its
rules is a mixfix production whose holes are typed by Lean's types and
grouped by precedence. A rendering is \emph{faithful} when it has exactly
one reading in that grammar, and that reading is the statement it was
rendered from, up to what the notation declares invisible: implicit
arguments, coercions, proofs, the names of hypotheses and bound variables,
and the premises and synonyms it names.

imprimatur checks this for each formula, by translation validation. As the
extractor renders a formula, it also builds the reading the formula should
have: an atom reads as itself, and a rule's rendering as the rule's pattern
with its holes' readings. It then parses the typeset formula, and
\verb|imprimatur check| reports, where the formula is, a formula with two
readings (showing both), with none, or with one that is not the expected
reading, a \emph{misread}. The last is what the check is for. A renderer
that left out the parentheses its holes need would print
$(\lambda r\, q.\, e) = M$ as $\lambda r\, q.\, e = M$, a formula that
reads back uniquely, and wrongly. Made to do so, imprimatur reports six
misreads in a twenty-page paper about another Lean development.

A few forms (a quantifier's binders, a case analysis, an associative chain)
are not given an expected reading, and only what is below them is
compared. There the check rests on the grammar mirroring the renderer, as
both are built from the same rules. In that twenty-page paper, 610 of the
662 formulas with one reading are compared in full, 21 in part, and 31 not
at all.

\paragraph{What the reader knows.} Reading back uses the types of the
formula's variables, which a reader has only if the text gives them. The
\verb|prose| view states the type of each variable it binds (\emph{Let
$b \in \mathit{Bundle}$}). A variable that a grammar or a figure leaves
free is known by its letter, as in BNF, where $p ::= \ldots$ says what $p$
ranges over. With \verb|latex_style categories := strict|, as in this
paper, imprimatur makes the convention exact: a variable whose type has a
metavariable is written with it and its own subscripts and primes, a
metavariable stands for one type, and two variables of one formula never
print alike. Each exception is an error. Appendix~\ref{app:notation} lists
the metavariables and the notation, both generated. A reader who knows
them has the grammar the formulas were checked in.

\paragraph{Together.} Theorem~\ref{thm:Soundness-for-blocks} and
translation validation make one claim. If the check reports nothing,
warnings included, then every block is the rendering of a trusted
constant's statement, and each formula in it reads back, in the paper's
notation, as that statement up to what the notation declares invisible.
The first half is proved for a model of the splicer; the second is checked,
for each formula of the paper, each time it is built.

\section{The trusted base}

What the two halves do not cover is trusted:
\begin{itemize}
\item Lean's kernel, and the extractor's reading of the environment.
\item The hash's resistance to collisions.
\item The parser that reads formulas back, and, where an expected reading
  is incomplete, the grammar's mirroring of the renderer.
\item The notation module, as the paper's dictionary. Lean checks its names,
  not its choice of symbols: a rule that wrote addition as $a - b$ would read
  back faithfully. Appendix~\ref{app:notation} puts the dictionary in the
  paper, where a reader can check it.
\item What the notation declares invisible. A hidden premise makes the
  theorem shown stronger than the one proved; a synonym is an equation the
  notation asserts and nothing checks.
\item The text of a block that is not a formula: the authors' text in
  slots, and the sentences of the \verb|prose| view (\emph{Let \ldots\
  Assume \ldots\ Then \ldots}), whose reading the view fixes.
\item \TeX{}. The checks stop at \TeX{} source, and two macros may print
  alike.
\item For attested claims, the person who checked the prose.
\end{itemize}

\begin{thebibliography}{1}
\bibitem{Pollack98}
R.~Pollack.
\newblock How to believe a machine-checked proof.
\newblock In \emph{Twenty-Five Years of Constructive Type Theory}. Oxford
  University Press, 1998.
\bibitem{Wiedijk12}
F.~Wiedijk.
\newblock Pollack-inconsistency.
\newblock \emph{Electronic Notes in Theoretical Computer Science}, 285,
  2012.
\end{thebibliography}

\appendix
\section{Notation}\label{app:notation}

The metavariables of this paper, and the notation it uses, generated from
its notation module. A variable with a subscript or a prime ranges over the
same type as its letter.

\begin{leanconventions}[sum=y0BasabyJh]{Conventions}{Metavariables and the types they range over}
\begin{longtable}{@{}l@{\quad}l@{\quad}p{.4\linewidth}@{\quad}l@{}}
\textit{Variable} & \textit{Type} & & \textit{Shown in}\\
\hline\endhead
$\sigma$ & $\mathit{Seg}$ & \leanslot{SpliceModel.Seg}{A segment of a generated block} & \textit{Seg}, p.~\pageref{gram:SpliceModel.Seg}\\
$p$ & $\mathit{Piece}$ & \leanslot{SpliceModel.Piece}{A piece of a document} & \textit{Piece}, p.~\pageref{gram:SpliceModel.Piece}\\
$k$ & $\mathit{Key}$ & \leanslot{SpliceModel.Key}{Keys: the names of Lean constants} & \textit{Piece}, p.~\pageref{gram:SpliceModel.Piece}\\
$x$ & $\mathit{Diag}$ & \leanslot{SpliceModel.Diag}{What \LeanCodeInline{check} reports} & \textit{Diag}, p.~\pageref{gram:SpliceModel.Diag}\\
$b$ & $\mathit{Bundle}$ & \leanslot{SpliceModel.Bundle}{A bundle: the entries of the keys a document uses} & \textit{Locality}, p.~\pageref{thm:Locality}\\
$d$ & $\mathit{Doc}$ & \leanslot{SpliceModel.Doc}{A document} & \textit{Locality}, p.~\pageref{thm:Locality}\\
\end{longtable}
A variable with a prime or a subscript ($\sigma'$, $\sigma_{1}$) ranges over the same type.
\end{leanconventions}

\begin{leannotation}[sum=wlP0ROlKQw]{Notation}{The notation of the paper}
\begin{longtable}{@{}l@{\quad}l@{\quad}p{.45\linewidth}@{}}
\multicolumn{3}{@{}l}{\textit{Lean}}\\
$\mathsf{List}(\alpha)$ & \texttt{List} & \leanslot{List/0}{Linked lists: ordered lists, in which each element has a reference to the next element}\\
$a \mathbin{::} l$ & \texttt{List.cons} & \leanslot{List.cons/0}{The list whose first element is \LeanCodeInline{head}, where \LeanCodeInline{tail} is the rest of the list}\\
$[]$ & \texttt{List.nil} & \leanslot{List.nil/0}{The empty list, usually written \LeanCodeInline{[]}}\\
$a \notin l$ & \texttt{Not} & \leanslot{Not/0}{\LeanCodeInline{Not p}, or \LeanCodeInline{¬p}, is the negation of \LeanCodeInline{p}}\\
$o \mathbin{\mathsf{else}} t$ & \texttt{Option.getD} & \leanslot{Option.getD/0}{Gets an optional value, returning a given default on \LeanCodeInline{none}}\\
$x$ & \texttt{Option.some} & \leanslot{Option.some/0}{Some value of type \LeanCodeInline{α}}\\
$\alpha \times \beta$ & \texttt{Prod} & \leanslot{Prod/0}{The product type, usually written \LeanCodeInline{α × β}}\\
$\mathit{String}$ & \texttt{String} & \leanslot{String/0}{A string is a sequence of Unicode scalar values}\\
\multicolumn{3}{@{}l}{\textit{SpliceModel}}\\
$\mathit{Body}$ & \texttt{Body} & \leanslot{SpliceModel.Body/0}{The body of a generated block}\\
$\mathsf{broken}(k)$ & \texttt{Diag.brokenRef} & \leanslot{SpliceModel.Diag.brokenRef/0}{the key has no entry}\\
$\mathsf{drift}(k)$ & \texttt{Diag.drift} & \leanslot{SpliceModel.Diag.drift/0}{the statement has changed since the prose was checked}\\
$\mathsf{stale}(k)$ & \texttt{Diag.stale} & \leanslot{SpliceModel.Diag.stale/0}{the block is not the key's rendering}\\
$\mathsf{untrusted}(k)$ & \texttt{Diag.untrusted} & \leanslot{SpliceModel.Diag.untrusted/0}{the key's proof is not trusted}\\
$\mathit{Env}(S)$ & \texttt{Env} & \leanslot{SpliceModel.Env/0}{A Lean environment: each key's statement, if it is a constant, and whether its proof is trusted}\\
$\mathit{stmt}_{E}(k)$ & \texttt{Env.stmt} & \leanslot{SpliceModel.Env.stmt/0}{the statement of each constant}\\
$\mathit{trusted}_{E}(k)$ & \texttt{Env.trusted} & \leanslot{SpliceModel.Env.trusted/0}{whether the proof of each constant is trusted}\\
$\mathsf{block}_{k}(\mathit{body})$ & \texttt{Piece.block} & \leanslot{SpliceModel.Piece.block/0}{a generated block for a key}\\
$\mathsf{checked}_{k}(h)$ & \texttt{Piece.check} & \leanslot{SpliceModel.Piece.check/0}{a check that hand-written prose matches the statement with this hash}\\
$\mathsf{text}(s)$ & \texttt{Piece.text} & \leanslot{SpliceModel.Piece.text/0}{hand-written text}\\
$\mathsf{gen}(s)$ & \texttt{Seg.gen} & \leanslot{SpliceModel.Seg.gen/0}{text generated from Lean}\\
$\mathsf{slot}_{i}(t)$ & \texttt{Seg.slot} & \leanslot{SpliceModel.Seg.slot/0}{a slot, whose text the authors own}\\
$\mathit{SlotId}$ & \texttt{SlotId} & \leanslot{SpliceModel.SlotId/0}{The identifiers of slots}\\
$\mathsf{wf}(b)$ & \texttt{WF} & \leanslot{SpliceModel.WF/0}{A bundle is well formed when no rendering has two slots with the same identifier}\\
$\mathsf{apply}_{b}(d)$ & \texttt{apply} & \leanslot{SpliceModel.apply/0}{\LeanCodeInline{imprimatur apply}: regenerate every block}\\
$\mathsf{check}_{b}(d)$ & \texttt{check} & \leanslot{SpliceModel.check/0}{\LeanCodeInline{imprimatur check}: what it reports about a document}\\
$\mathsf{export}(E)$ & \texttt{extract} & \leanslot{SpliceModel.extract/0}{\LeanCodeInline{imprimatur-export}: the bundle of an environment, given how statements are rendered and hashed}\\
$\mathsf{fill}(\mathit{old}, r)$ & \texttt{fill} & \leanslot{SpliceModel.fill/0}{Fill a rendering's slots: each keeps the authors' text, if they wrote some, else the rendering's}\\
$\mathsf{ids}(r)$ & \texttt{ids} & \leanslot{SpliceModel.ids/0}{The identifiers of a body's slots}\\
$l[i]$ & \texttt{lookup} & \leanslot{SpliceModel.lookup/0}{The text of slot \LeanCodeInline{i}, the first with that identifier}\\
$\lfloor d \rfloor$ & \texttt{outline} & \leanslot{SpliceModel.outline/0}{The outline of a piece: a block without its body}\\
$\mathsf{slots}(r)$ & \texttt{slots} & \leanslot{SpliceModel.slots/0}{The slots of a body, with their text}\\
\end{longtable}
\end{leannotation}

\end{document}
